\documentclass[11pt]{article}

\usepackage[margin=1in]{geometry}
\usepackage{amsmath,amssymb}
\usepackage{microtype}
\usepackage{newtxtext,newtxmath}
\usepackage{parskip}
\usepackage{enumitem}
\usepackage[round,authoryear]{natbib}
\usepackage[
  colorlinks=true,
  linkcolor=black,
  citecolor=black,
  urlcolor=blue
]{hyperref}

\setlength{\bibsep}{5pt}
\setlist{nosep}

\newcommand{\term}[1]{\emph{#1}}
\newcommand{\modality}{\mathcal{M}}
\newcommand{\contentobj}{\mathcal{C}}
\newcommand{\receivedobj}{\mathcal{R}}

\title{\textbf{Modality-Heading}\\
\large Communication, Cognitive Operators, and Dependence on Form}

\author{Flyxion\\
\small Independent Researcher}

\date{September 2026}

\begin{document}

\maketitle

\begin{abstract}
Communication is commonly modeled as the transmission of content through
channels that alter its presentation without materially altering what is
communicated. This model is useful only when the effects of mediation can be
ignored. Prose, mathematical notation, typography, comics, music, audiovisual
narrative, and conversation organize attention, temporality, affect,
inference, participation, and memory in different ways. A modality is
therefore not merely a carrier of communicative material but an operator that
helps determine the object ultimately received.

This essay introduces \term{modality-heading}, the progressive attachment of
cognitive reward, perceived intelligibility, and epistemic acceptance to a
preferred transformation of communicative material. The concept generalizes
\term{typefaceheading}, in which inscriptional form produces an affect of
semantic significance that may exceed the warrant supplied by the inscription's
content. Modality-heading differs from conventional accounts of media
addiction, algorithmic enclosure, and ideological echo chambers. Its central
problem is not simply what information is selected or how compulsively a
medium is used, but what must happen to information before it can be admitted
as meaningful.

Drawing upon transitional-object theory, surrogate-attachment research,
wireheading, media theory, and human--computer interaction, the essay
distinguishes five functions of modality: expression, regulation, refuge,
constitution, and enclosure. A modality may communicate an experience,
regulate its intensity, make otherwise intolerable material approachable, or
supply a performative scaffold through which an agent becomes capable of
action. These functions become enclosing when the preferred form loses
corrective permeability and begins excluding whatever cannot be assimilated
to its grammar.

Adaptive conversational systems intensify this possibility because they
operate as personalized modality converters. They can transform technical
arguments, emotional experiences, narratives, and formal structures into a
responsive conversational register while attenuating many of the reciprocal
demands imposed by human interlocutors. The resulting danger is not merely
mistaken anthropomorphism or excessive attachment. It is the emergence of
format-conditioned admissibility: a communicative environment in which the
world remains available only insofar as it can be reformatted into a preferred
mode of reception. The essay concludes by proposing cross-modal literacy and
corrective permeability as criteria for distinguishing functional mediation
from modal enclosure.
\end{abstract}

\newpage

\section{The Non-Neutrality of Form}

A familiar model of communication begins with a message, places that message
inside a channel, and asks how accurately it arrives at its destination. The
channel may introduce noise, reduce fidelity, or impose limits upon
transmission, but the message itself remains conceptually prior to the means
by which it travels. This separation is indispensable in information theory,
where the technical problem is to reproduce a selected message at another
point despite uncertainty introduced during transmission
\citep{shannon1948}. It becomes less satisfactory when extended from signal
reproduction to human understanding. A sentence can be transmitted without
error while being misunderstood, dismissed, aestheticized, trivialized, or
made newly persuasive by the form in which it appears.

The inadequacy does not arise because human communication is uniquely
irrational. It arises because communicative forms do more than preserve or
degrade signals. They organize the conditions under which signals become
available for cognition. Written prose permits rereading, qualification,
recursive embedding, and retrospective comparison. Conversation introduces
turn order, interruption, repair, hesitation, and the possibility of an
immediate reply. Mathematical notation makes some relations inspectable at a
glance while excluding much of the contextual information carried by ordinary
language. Music organizes expectation through duration, repetition, tension,
and release. Comics distribute temporal sequence across space and permit the
reader to vary the duration of inspection without changing the represented
order of events. These differences are not ornaments applied after meaning
has been constituted. They participate in its constitution.

The claim that form participates in meaning is easily weakened into the
truism that presentation matters. The stronger claim developed here is that a
modality operates upon communicative material. It selects relations, assigns
salience, controls pacing, distributes agency, regulates affective intensity,
and establishes characteristic forms of completion. A modality does not
merely make an independently complete object easier or harder to apprehend. It
helps determine what kind of object can be apprehended at all.

Let \(\contentobj\) denote communicative material, \(\modality\) a modality,
\(a\) a receiving agent, and \(e\) the environment in which reception occurs.
The received object \(\receivedobj\) may then be represented schematically as

\begin{equation}
  \receivedobj = \modality(\contentobj,a,e).
  \label{eq:operator-model}
\end{equation}

Equation~\ref{eq:operator-model} is not intended as a quantitative model. Its
purpose is to prevent the received object from being identified automatically
with the material supplied to the medium. Reception depends upon an operation
whose result varies with the capacities, expectations, history, and situation
of the recipient. The same proof encountered by a trained mathematician, a
novice, and an automated theorem checker does not produce the same received
object, even when the marks presented to each are identical. Likewise, the
same words spoken during an argument, printed in a textbook, sung in a refrain,
or displayed as an institutional warning enter different structures of
attention and action.

The term \term{modality} is used here in a deliberately broad sense. It
includes sensory and representational forms, such as speech, writing, image,
music, and mathematical notation, but it also includes organized genres and
interfaces when these impose sufficiently stable operations upon reception.
A detective narrative distributes evidence differently from a textbook. A
melodrama makes interpersonal causation unusually conspicuous. A search
interface presents knowledge as retrieval from a field of alternatives,
whereas a conversational interface presents it as the unfolding response of
an addressable interlocutor. These cases differ in scale, but each establishes
constraints upon what can appear, in what order, and with what presumptive
significance.

No claim of absolute determination follows. A medium does not compel every
recipient to derive the same interpretation, nor does genre eliminate
individual variation. The operator model instead asserts that interpretation
occurs within a transformed field. An agent may resist, reinterpret, parody,
or misuse a modality, but such resistance still takes place against the
affordances and expectations the modality establishes. A reader may refuse
the suspense of a mystery by reading its final page first, yet this act has
its significance only because the form ordinarily controls the sequence in
which information is disclosed.

This account also differs from a strong doctrine according to which the
medium wholly determines the message. The frequently cited proposition that
``the medium is the message'' correctly directs attention toward changes in
scale, pace, and social organization produced by media themselves
\citep{mcluhan1964}. The present argument is narrower. It does not identify
medium and message, nor does it deny the persistence of claims across
representations. It asks instead which properties survive a transformation,
which are introduced by it, and which disappear because the receiving
modality has no corresponding operation.

Suppose that a communicative object is translated from modality
\(\modality_1\) into modality \(\modality_2\). It is tempting to describe the
translation as successful whenever both versions concern the same subject.
Subject continuity, however, is a weak criterion. A prose explanation of a
musical composition may accurately identify its key, tempo, historical
setting, and formal divisions while preserving almost none of the temporal
experience through which those features become musically intelligible.
Conversely, a performance may preserve the operative musical relations while
communicating none of the historical claims contained in the prose account.
Neither representation is simply defective. Each preserves a different set
of invariants.

Cross-modal identity must therefore be indexed to a preservation criterion.
If \(P\) denotes a property or relation intended to survive translation, then
a minimal claim of continuity can be written as

\begin{equation}
  P\!\left(\modality_1(\contentobj)\right)
  \cong
  P\!\left(\modality_2(\contentobj)\right),
  \label{eq:modal-invariant}
\end{equation}

where \(\cong\) denotes an equivalence appropriate to the property under
consideration rather than literal identity. The preserved property might be a
proposition, inferential dependency, emotional orientation, temporal order,
procedural constraint, character relation, or rhythmic structure. Without
specifying \(P\), the assertion that two modalities contain the same
information conceals the standard by which sameness has been judged.

The concealment is consequential because the ease with which a property
appears in one modality can be mistaken for a property of the material
itself. A diagram may make a relation immediately visible, encouraging the
belief that the relation is intrinsically simple. Formal notation may make an
inference compact, encouraging the belief that its premises are
uncontroversial. Narrative may make a causal sequence memorable, encouraging
the belief that its events possess a single intelligible direction. A
conversational explanation may feel responsive and complete because it
resolves uncertainty at the pace at which uncertainty is expressed. In each
case, an affordance of the modality can be misrecognized as evidence about
the represented object.

The distinction can be stated as follows:

\begin{equation}
  \text{ease of reception}
  \not\Rightarrow
  \text{adequacy of representation}.
  \label{eq:ease-adequacy}
\end{equation}

The converse is equally important. Difficulty does not establish adequacy.
An obscure presentation is not made rigorous merely by resisting
comprehension, and a familiar form is not epistemically compromised merely
because it is pleasurable. The problem is not ease itself but the failure to
attribute ease to its sources. A competent analysis must separate the
properties of the represented material from the operations through which
those properties become available.

This separation provides the starting point for the theory of
modality-heading. If modalities produce distinctive cognitive and affective
rewards, agents may learn to seek not only particular messages but particular
transformations of messages. A preferred form can then move from convenience
to expectation, from expectation to dependency, and from dependency to an
implicit criterion of admissibility. Before developing that general process,
however, it is useful to examine a restricted case in which the operation of
form is unusually visible: the capacity of typography to manufacture an
affect of meaning before the represented claim has been assessed.

\section{Typefaceheading}

Typography offers a restricted case in which the operation of modality is
difficult to overlook. The lexical content of an inscription may remain
stable while changes in typeface, spacing, weight, alignment, color, motion,
or graphical corruption alter its apparent voice. The same sentence can
appear juridical, intimate, mechanical, archaic, commercial, devotional, or
unstable without any substitution of words. These changes do not necessarily
modify the proposition expressed, but they modify the conditions under which
the proposition is encountered.

This effect is sometimes dismissed as decoration, as though typography were
a detachable surface laid over an otherwise complete linguistic object.
Such a description mistakes lexical recoverability for communicative
identity. Written language is never encountered without material form.
Typeface, scale, spacing, contrast, and arrangement govern segmentation,
hierarchy, emphasis, continuity, and anticipated genre. Even a conventionally
``neutral'' typeface communicates through accumulated conventions of
institutional use. Neutrality is not the absence of typographic operation but
the successful naturalization of one typographic regime.

The operation becomes conspicuous when those conventions are violated.
Consider an ordinary phrase rendered with dense stacks of combining Unicode
characters. Its letters may remain recoverable, yet the inscription appears
infected, haunted, damaged, computationally unstable, or insufficiently
contained by the line on which it is printed. The semantic effect arises
partly from a conflict between lexical persistence and graphical
disintegration. The reader recognizes the phrase while simultaneously
encountering an inscription that seems to exceed or corrupt its own encoding.

The marks do not establish that the phrase refers to instability, corruption,
or haunting. They produce those associations through the behavior of the
inscription itself. Form becomes locally demonstrative: the text appears to
undergo what it might otherwise only describe. This performative coincidence
can be rhetorically powerful, but it can also conceal the difference between
depicting a condition and establishing a claim about it.

The term \term{typefaceheading} denotes the acquisition of cognitive or
affective reward through the inscriptional treatment of language. The
compound is modeled upon \term{wireheading}, but the mechanism is less direct.
The reader still encounters lexical and syntactic material. What changes is
the degree to which the anticipated reward of semantic achievement can be
elicited by graphical form before, or in the absence of, corresponding
conceptual work.

Typefaceheading should not be identified with typographic pleasure in
general. Legibility, beauty, visual rhythm, and stylistic coherence are
legitimate dimensions of written communication. Nor does the term imply that
typography is external to meaning. A poem arranged spatially on a page, a
mathematical proof divided into inferential units, and a warning displayed in
large capital letters cannot be fully described by recording their lexical
content alone. The distinction concerns the relation between the effect
produced by form and the warrant supplied by what has been expressed.

Let \(S\) denote semantic material and \(T\) its typographic treatment. The
experienced significance of an inscription may be represented schematically
as

\begin{equation}
  \Sigma = f(S,T,a,e),
  \label{eq:typographic-significance}
\end{equation}

where \(a\) and \(e\) again represent the receiving agent and environment.
The equation makes no claim that semantic and typographic contributions can
be measured independently. It asserts only that experienced significance
cannot be attributed to \(S\) alone.

The relevant epistemic caution is

\begin{equation}
  \text{semantic affect}
  \not\Rightarrow
  \text{semantic warrant}.
  \label{eq:affect-warrant}
\end{equation}

An inscription may feel profound, dangerous, ancient, technical, or
revelatory without providing reasons adequate to those impressions.
Conversely, an austere or familiar presentation may contain an argument of
considerable consequence without producing an immediate affect of
significance. The distinction is elementary when stated explicitly, yet
interfaces routinely profit from its concealment.

Typographic form can imitate the visible residues of difficulty. Dense
notation, specialist symbols, unusual scripts, marginal annotations, damaged
characters, and elaborate diagrams may suggest that an inscription has passed
through processes of compression, transmission, scholarship, secrecy, or
technical refinement. In some cases it has. In others, the residues are
introduced directly as stylistic signals. The appearance of having survived a
history substitutes for the history itself.

This substitution is especially effective when readers have learned reliable
associations between graphical form and institutional practice. Monospaced
text evokes programming and terminal interaction. Blackletter evokes
historical authority, nationalism, liturgy, or antiquarianism according to
context. Geometric sans serif type may evoke administrative modernity or
technical cleanliness. Handwriting can suggest intimacy and singular
presence even when mechanically reproduced. Typographic conventions become
compressed social histories that can be activated without reproducing the
conditions in which their associations developed.

A useful distinction can therefore be made between \term{inscriptional
evidence} and \term{inscriptional simulation}. Inscriptional evidence consists
of formal features generated by the process through which a text was
produced, transmitted, or used. Corrections in a manuscript may record
revision; inconsistent type in an archival document may indicate composite
production; compression artifacts may reveal a history of digital
transmission. Inscriptional simulation introduces comparable features in
order to evoke such a history without being caused by it. The distinction
does not make simulation illegitimate. Fiction, design, parody, and
pedagogy routinely depend upon it. It becomes epistemically relevant when
the simulated residue is treated as evidence of the process it merely
resembles.

Typefaceheading emerges through repeated reinforcement of such effects. A
reader encounters inscriptions whose visual treatment reliably precedes
experiences of novelty, expertise, transgression, intimacy, or discovery.
The treatment itself then acquires anticipatory force. It signals not only
what kind of text is present but what kind of cognitive reward should be
expected from it. Eventually, a familiar graphical signature can deliver part
of that reward independently of the text's success.

This process can be described as a migration of valuation. Initially,
typographic form directs attention toward semantic or practical features.
Subsequently, the form becomes associated with the reward produced by those
features. Finally, reproduction of the form elicits the reward with a reduced
requirement for semantic achievement:

\begin{equation}
  T \rightarrow S \rightarrow \Sigma
  \qquad\Longrightarrow\qquad
  T \rightarrow \Sigma.
  \label{eq:typographic-short-circuit}
\end{equation}

The second relation does not imply that semantic material disappears. It
indicates that the path through it has ceased to be necessary for some part
of the response. The typographic surface has become a shortened route to the
affect of significance.

The word \term{heading} is therefore appropriate in two senses. It names an
orientation toward a preferred source of stimulation, as in wireheading, and
it recalls the typographic heading as a device that assigns importance before
the subordinate text has been read. A heading organizes expectation. Under
typefaceheading, the organization of expectation becomes partly sufficient
to produce the expected valuation.

The phenomenon is easiest to see in extreme inscriptions, but it is not
confined to them. Academic formatting, interface design, branding, document
templates, citation density, and mathematical display can all participate in
the same transfer. A displayed equation may be indispensable to an argument,
or it may function chiefly as a visual certificate that an argument belongs
to a rigorous class. The difference cannot be determined from the presence
of notation alone. It depends upon whether the notation performs inferential
work that prose cannot perform as precisely, and whether the surrounding
claims remain answerable to it.

Typefaceheading is thus not a denunciation of stylization. It is a diagnostic
concept for cases in which inscriptional form supplies a reward normally
attributed to semantic accomplishment. Its value lies in making a small
mechanism visible. Once that mechanism is recognized, there is no reason to
restrict it to typography. Any stable communicative modality can become
associated with characteristic rewards, and any such association can alter
the conditions under which subsequent material is admitted as intelligible.
Typefaceheading is therefore a local instance of a more general dependence
upon form.

\section{From Typefaceheading to Modality-Heading}

The mechanism described in the preceding section is not specific to written
form. Typography makes it unusually visible because lexical material can
remain approximately constant while inscriptional treatment changes around
it. Comparable transfers occur whenever a communicative form becomes
reliably associated with a characteristic cognitive reward. Narrative can
supply causal closure; mathematics can supply relational compression;
conversation can supply responsiveness; music can supply patterned affect;
visualization can supply apparent simultaneity; and interactive media can
supply agency. In each case, the reward may migrate from what the modality
helps disclose toward the operation of the modality itself.

The term \term{modality-heading} denotes this more general process:

\begin{quote}
\textbf{Modality-heading} is the progressive attachment of cognitive reward,
perceived intelligibility, and epistemic acceptance to a preferred
transformation of communicative material.
\end{quote}

The definition contains three distinguishable components. \term{Cognitive
reward} concerns the felt satisfaction associated with comprehension,
recognition, completion, surprise, fluency, or successful prediction.
\term{Perceived intelligibility} concerns whether material appears organized
enough to be understood. \term{Epistemic acceptance} concerns whether the
material is treated as credible, serious, relevant, or worthy of further
attention. These components frequently reinforce one another, but they need
not coincide. A presentation may be pleasurable without appearing credible,
or credible without being immediately intelligible. Modality-heading becomes
especially powerful when a single form coordinates all three.

A preference for one modality over another does not by itself constitute
modality-heading. Preferences reflect training, aptitude, accessibility,
purpose, and context. A mathematician may reasonably prefer notation when
relations would become cumbersome in prose. A musician may apprehend a
harmonic transition more precisely by hearing it than by reading a verbal
description. A reader with limited vision may prefer speech synthesis, while
a reader who needs to compare distant passages may prefer a stable written
page. These are functional selections among unequal affordances.

The relevant transition occurs when preference becomes a condition of
admissibility. Material presented outside the favored modality no longer
appears merely inconvenient. It appears deficient in meaning, rigor,
emotional reality, or truth. The receiving agent begins to evaluate the
material according to whether it has undergone the expected transformation.

Let \(A(\contentobj,\modality,a)\) denote the degree to which communicative
material \(\contentobj\), presented through modality \(\modality\), is admitted
by agent \(a\) as worthy of cognitive uptake. A strong modal dependence exists
when admissibility varies more with the modality than with the relevant
properties of the material:

\begin{equation}
  \left|
  A(\contentobj,\modality_1,a)
  -
  A(\contentobj,\modality_2,a)
  \right|
  \gg
  \Delta_P(\contentobj),
  \label{eq:modal-dependence}
\end{equation}

where \(\Delta_P(\contentobj)\) represents differences relevant to the
declared epistemic purpose \(P\). This expression is schematic rather than
metric. It identifies the condition in which a change of form produces a
larger difference in acceptance than any corresponding change in the
properties that acceptance is supposed to track.

Modality-heading can therefore be understood as a displacement of the
admissibility function. An agent initially uses a modality to gain access to
properties of an object. Through repetition, the modality becomes associated
with successful access. The presence of the modality is then treated as
evidence that access has occurred:

\begin{equation}
  \text{access through }\modality
  \longrightarrow
  \text{successful understanding}
  \longrightarrow
  \modality\text{ as evidence of understanding}.
  \label{eq:admissibility-displacement}
\end{equation}

The displacement resembles typefaceheading but operates at a larger scale. A
person may learn to regard narrative coherence as evidence of causal
adequacy, mathematical compactness as evidence of ontological depth,
conversational fluency as evidence of comprehension, visual immediacy as
evidence of transparency, or emotional intensity as evidence of importance.
None of these associations is wholly arbitrary. Narrative, notation,
dialogue, visualization, and affect can genuinely disclose relevant
properties. The error consists in converting a defeasible affordance into an
automatic certificate.

\subsection{Distinction from Media Addiction}

Modality-heading overlaps with media addiction but is not reducible to it.
Accounts of behavioral addiction usually emphasize compulsive repetition,
impaired control, escalating use, withdrawal-like distress, and continuation
despite adverse consequences. These features concern the regulation of
behavior around a medium. Modality-heading concerns the regulation of
intelligibility through a medium.

A person may spend very little time with a modality while nevertheless
granting it disproportionate epistemic authority. A short mathematical
display may be treated as more serious than several pages of verbal argument
even when the display merely restates an unsupported premise. A brief video
may be treated as decisive because it produces perceptual immediacy, while a
careful statistical analysis is dismissed as abstract. Conversely, a person
may use a medium intensively without requiring all other material to assume
its form.

The difference can be expressed by separating frequency of use \(F\) from
modal dependence \(D\):

\begin{equation}
  F(\modality,a)
  \not\Rightarrow
  D(\modality,a),
  \qquad
  D(\modality,a)
  \not\Rightarrow
  F(\modality,a).
  \label{eq:frequency-dependence}
\end{equation}

High frequency may contribute to dependence through reinforcement, but the two
variables identify different phenomena. Media addiction asks whether use has
become behaviorally compulsive. Modality-heading asks whether a form has
become epistemically compulsory.

The distinction also separates modality-heading from an ideological echo
chamber. An echo chamber restricts exposure to claims, sources, or social
positions that reinforce an existing view. Modality-heading may persist even
when the content is heterogeneous and adversarial. An agent can encounter
opposing political positions, unfamiliar scientific claims, and diverse
cultural material while requiring all of them to arrive as short videos,
formal models, serialized narratives, or conversational summaries. Content
diversity does not guarantee modal diversity.

Algorithmic filtering introduces another related but distinct process.
Recommendation systems determine which items become available or salient.
Modality-heading concerns what happens to an item once selected and the form
it must assume before it can be accepted. Selection and transformation can
interact: an algorithm may favor material already optimized for a preferred
modality, while an adaptive interface may transform otherwise unsuitable
material into that modality. Nevertheless, the analytical distinction should
be retained:

\begin{equation}
  \text{selection determines what arrives},
  \qquad
  \text{modality determines how it can arrive}.
  \label{eq:selection-transformation}
\end{equation}

\subsection{Wireheading and the External Remainder}

The suffix \term{-heading} invites comparison with wireheading, in which a
reward system is stimulated without the worldly activity or condition that
would ordinarily produce the reward. An organism need not obtain nourishment,
security, achievement, or social recognition if the reinforcing signal can
be supplied more directly. The reward becomes detached from the state of
affairs it once helped evaluate.

Typefaceheading shortens the path to semantic affect by reproducing visible
cues associated with significance. Modality-heading performs a related
shortening, but it does not ordinarily bypass external material altogether.
The world continues to provide events, claims, evidence, and resistance.
These inputs are reformatted through a preferred operator before they are
admitted.

The contrast can be stated compactly:

\begin{align}
  \text{wireheading} &: \quad
  \text{worldly condition} \;\;\text{bypassed},
  \label{eq:wireheading}\\
  \text{modality-heading} &: \quad
  \text{worldly material} \;\;\text{reformatted}.
  \label{eq:modality-heading}
\end{align}

The presence of an external remainder prevents modality-heading from being
treated as pure hallucination or direct stimulation. A narrative may convey
historical evidence; a visualization may reveal a genuine distribution; a
conversation may clarify a real argument. The concern is not that nothing
from the world survives. It is that the preferred transformation increasingly
determines which parts survive and what status they acquire.

This distinction also explains why modality-heading can remain productive for
long periods. A preferred form often develops because it works. It reduces
complexity, coordinates attention, permits action, regulates distress, or
makes otherwise inaccessible structures available. Its success supplies the
reinforcement from which dependency grows. Enclosure is therefore rarely
introduced as enclosure. It begins as an effective solution to a real problem
of access.

\subsection{The Gradient from Preference to Enclosure}

The development of modality-heading can be represented as a progression:

\begin{equation}
  \text{affordance}
  \longrightarrow
  \text{preference}
  \longrightarrow
  \text{expectation}
  \longrightarrow
  \text{requirement}
  \longrightarrow
  \text{enclosure}.
  \label{eq:heading-gradient}
\end{equation}

At the level of affordance, a modality makes some operation easier. At the
level of preference, the agent selects it when alternatives remain usable.
At the level of expectation, its absence produces friction or
disappointment. At the level of requirement, material outside the modality
fails to receive adequate cognitive uptake. At the level of enclosure, the
agent loses the ability to recognize the modality as one transformation among
others.

The final stage is not defined by intense enjoyment. It is defined by the
disappearance of modal contingency. The preferred form no longer appears to
be a form. It appears to be the natural shape of intelligibility itself.

This gradient is diagnostic rather than strictly developmental. An agent may
occupy different positions for different domains, move backward through
training, or deliberately tolerate dependence when accessibility requires
it. Nor is maximal modal flexibility always desirable. Expertise depends
partly upon specialized representational systems whose power cannot be
reproduced in every alternative form. The relevant question is whether the
specialized modality remains answerable to objects, consequences, and
representations beyond its own grammar.

To answer that question, modality must be examined not only as a device for
presenting information but as a means of regulating contact with what exceeds
the agent. The next section therefore turns from admissibility to mediation.
Transitional-object theory provides a vocabulary for understanding how forms
can make difficult realities approachable without yet becoming substitutes
for reality itself.

\section{Transitional Objects and Transitional Media}

The preceding account describes how a useful communicative form can become an
implicit condition of intelligibility. It does not yet explain why particular
forms acquire such regulatory importance. Transitional-object theory offers
one point of entry because it treats mediation not as an accidental barrier
between subject and world but as a necessary region in which their relation is
gradually organized.

Winnicott introduced the transitional object to describe an infant's
attachment to objects such as blankets, pieces of cloth, or soft toys
\citep{winnicott1953}. Such an object is neither experienced simply as part of
the self nor recognized in the same manner as an independently existing
person. It occupies an intermediate domain between subjective creation and
external reality. Its theoretical importance lies less in its material
composition than in the ambiguity it sustains.

This ambiguity should not be mistaken for a defective classification awaiting
correction. The transitional object is useful because the question of whether
it was created or discovered need not be settled immediately. It provides a
stable surface across which dependence, absence, imagination, and externality
can be negotiated. The object does not merely represent comfort. Its
availability, texture, portability, and persistence participate in the
regulation of experience.

The adjective \term{transitional} can misleadingly suggest that the object is
important only because it will eventually be discarded. What matters is not
the disappearance of mediation but the kind of movement mediation permits.
Language, ritual, art, technical notation, and shared conventions continue to
occupy intermediate regions throughout adult life. They connect private
experience with publicly available forms without eliminating the difference
between them. The developmental question is therefore not whether transitional
structures survive, but whether they continue to support movement beyond
themselves.

This yields a distinction central to the present argument:

\begin{equation}
  \text{transitional support}
  \not\equiv
  \text{permanent enclosure}.
  \label{eq:transition-enclosure}
\end{equation}

A support is transitional when it increases the range of relations an agent
can sustain. It becomes enclosing when continued regulation depends upon
preventing contact with whatever exceeds the support's organizing grammar.
The same object or practice may perform either function at different times.
No intrinsic property of softness, repetition, familiarity, or responsiveness
settles the matter.

The concept can be extended cautiously from objects to media. A
\term{transitional medium} is a representational or interactive form that
regulates the distance between an agent and material that cannot yet be
encountered directly, continuously, or in its full complexity. Narrative can
make diffuse historical forces approachable through agents and events.
Mathematical notation can stabilize relations that exceed unaided working
memory. Music can organize affect that resists propositional articulation.
Dialogue can externalize an unfinished line of reasoning so that it becomes
available for inspection and revision.

The phrase does not imply that the mediated object exists in a purer form
somewhere beyond representation. Direct access is not always possible, and in
many domains it is not conceptually coherent. A mathematical structure is not
made less mediated by being represented in notation, nor is an emotion
necessarily more authentic when it remains unnamed. Transitional mediation
should instead be evaluated by the relations it enables and the differences it
preserves.

\subsection{The Decomposition of Relational Functions}

Harlow's surrogate-mother experiments provide a severe example of relational
functions being separated and assigned to engineered objects. Infant rhesus
macaques were presented with surrogate figures that distributed nourishment
and tactile contact across different apparatuses. Their preference for a soft
cloth-covered surrogate demonstrated that attachment behavior could not be
explained by feeding alone \citep{harlow1958}.

The experiments are frequently summarized as a discovery that comfort matters
more than nourishment. That summary is too simple. Their broader implication
is that a relationship ordinarily encountered as a practical unity can be
decomposed into functions. Nourishment, warmth, tactile contact, security,
responsiveness, familiarity, and mobility need not be delivered by the same
entity. Once separated, these functions can be varied, withheld, recombined,
and measured.

The decomposition is analytically powerful and ethically compromised. The
experiments did not passively observe naturally occurring attachment. They
manufactured deprivation, restricted ordinary social development, and then
measured adaptation to impoverished alternatives. The resulting behavior
cannot be interpreted independently of the conditions that produced the need
for the surrogate. An engineered substitute may appear unusually effective
when the experiment has first removed competing sources of attachment.

This caveat applies beyond animal research. A technological system can appear
to satisfy a human need with exceptional efficiency when the surrounding
environment has made ordinary satisfaction expensive, unreliable, or
unavailable. The success of the substitute cannot be evaluated solely by
measuring attachment to it. One must also ask which relations were weakened,
withheld, or reorganized before the substitute became comparatively
attractive.

Nevertheless, the effects supplied by a surrogate are not unreal merely
because the surrogate is artificial. Tactile regulation produced by a cloth
surface is still tactile regulation. Predictability produced by a repeated
interface is still predictability. A responsive system may help organize
attention or reduce distress even if it does not reproduce the living
relationship from which those effects were abstracted.

The distinction is between an effect and the relational totality with which
that effect is ordinarily associated:

\begin{equation}
  \text{reproduction of a relational effect}
  \not\Rightarrow
  \text{reproduction of the relationship}.
  \label{eq:effect-relationship}
\end{equation}

The converse error should also be avoided:

\begin{equation}
  \text{absence of the original relationship}
  \not\Rightarrow
  \text{absence of a real effect}.
  \label{eq:relationship-effect}
\end{equation}

Together, Equations~\ref{eq:effect-relationship}
and~\ref{eq:relationship-effect} prevent two symmetrical reductions. The first
mistakes successful simulation of a function for recreation of the whole. The
second dismisses every benefit of mediation because the original relation is
absent. Both obscure the actual organization of the interaction.

\subsection{Protective Modalities}

A modality can become transitional by reducing the intensity at which
material is encountered. Familiar genres are especially capable of this
operation because they supply stable expectations about pacing, causality,
emotional range, and resolution. Material that would otherwise be
disorganizing can be admitted after conversion into a form whose rules are
already known.

The conversion does not necessarily erase the original material. It may
preserve recognizable relations while changing their affective availability.
Threat can become suspense, grief can become melody, conflict can become
comedy, and uncertainty can become a problem with a determinate solution
space. The resulting representation is neither an intact copy nor a complete
avoidance. It is a regulated encounter.

An unusually explicit fictional example occurs in \emph{Mr.\ Robot}, where
intolerable experience is reorganized through the conventions of a syndicated
family sitcom. Domestic staging, canned laughter, exaggerated performance,
and familiar television rhythm do not merely comment upon the character's
condition. They form a temporary protective environment through which the
condition can be represented without being encountered at full intensity.
The sitcom modality operates as psychic insulation.

The example matters because it shows that modal transformation can be
adaptive even when it is manifestly unreal. The protective form permits
partial contact where direct contact would be overwhelming. Its distortion is
not simply an epistemic defect; it is part of the mechanism by which
representation remains possible.

Let \(I(\contentobj,a)\) denote the unregulated intensity of material
\(\contentobj\) for agent \(a\), and let
\(I_{\modality}(\contentobj,a)\) denote its intensity after modal
transformation. A protective modality produces a relation of the form

\begin{equation}
  I_{\modality}(\contentobj,a)
  <
  I(\contentobj,a),
  \qquad
  P\!\left(\modality(\contentobj)\right)
  \cong
  P(\contentobj),
  \label{eq:protective-modality}
\end{equation}

where \(P\) identifies whatever relation must be preserved for the encounter
to remain connected to its object. The modality reduces intensity while
retaining a relevant invariant.

The preservation condition is essential. If no significant relation survives,
the modality functions less as regulated contact than as displacement. Yet
the amount that must survive cannot be specified independently of the agent's
capacity. A minimal invariant may be sufficient at one stage and evasive at
another. Transitional value is therefore temporally indexed.

\subsection{From Regulation to Refuge}

Protective mediation can be described through a sequence:

\begin{equation}
  \text{unmanageable material}
  \longrightarrow
  \text{modal transformation}
  \longrightarrow
  \text{regulated contact}
  \longrightarrow
  \text{possible integration}.
  \label{eq:regulated-contact}
\end{equation}

The sequence remains transitional only if the final term remains possible.
When regulation becomes an end in itself, the medium changes from a bridge
into a refuge. Refuge is not inherently pathological. Withdrawal can preserve
an agent whose immediate environment exceeds available capacities. Repetition,
familiarity, and controlled distance may be prerequisites for later
engagement.

The difficulty arises when the refuge reorganizes the agent's expectations so
thoroughly that unregulated contact becomes progressively less tolerable. The
medium initially selected because it reduces overwhelming intensity can
reduce the capacity to encounter intensity outside that medium. Protection
then produces the vulnerability that appears to justify further protection.

This feedback relation can be represented schematically as

\begin{equation}
  \text{relief through }\modality
  \longrightarrow
  \text{reduced exposure outside }\modality
  \longrightarrow
  \text{reduced extra-modal tolerance}
  \longrightarrow
  \text{greater dependence on }\modality.
  \label{eq:protective-feedback}
\end{equation}

The cycle is not inevitable. A protective modality may instead increase
capacity by permitting graded exposure, reflective distance, rehearsal, and
comparison. The difference lies in whether regulation returns the agent to a
wider field of action or continually narrows the conditions under which action
remains possible.

A transitional medium should therefore be evaluated by its effects upon
subsequent reachability. If its use enlarges the set of states, relations, and
representations the agent can approach, it functions as scaffolding. If it
makes continued functioning conditional upon repeated return to the same form,
it approaches enclosure.

The analysis so far has treated modalities as means of expressing or
regulating material that already presses upon an agent. Some modalities do
more. They provide scripts, roles, gestures, and expectations through which an
agent becomes capable of actions that were not previously organized. In such
cases, mediation is not merely protective. It is constitutive.

\section{Constitutive Modalities and Functional Fiction}

A protective modality regulates contact with material that would otherwise
exceed an agent's present capacity. A constitutive modality performs a
different operation. It supplies a form through which capacities become
coordinated into action. Its contribution is not limited to representing or
attenuating an antecedent state. By providing roles, scripts, gestures,
expectations, and criteria of completion, it helps produce the agent who acts.

This constitutive function is familiar in institutional life. A uniform does
not merely communicate an occupational identity to observers. It changes the
wearer's posture, obligations, permissible actions, and relation to a shared
procedure. Technical notation does not simply record thoughts already
completed in ordinary language. It makes sequences of transformation
available that could not be maintained reliably without inscription. Ritual
does not merely describe commitment; it coordinates bodies, attention, and
memory so that commitment can acquire a durable public form.

The same structure appears in self-narrative. An agent interprets present
conditions through an account of who they are, what kind of situation they
occupy, and which actions are available to a person of that kind. Such an
account may be descriptively imperfect while remaining organizationally
effective. Indeed, some self-narratives derive their efficacy from projecting
a coherence or capability that the agent does not yet possess.

This possibility requires a distinction between representational accuracy and
organizational efficacy. Let \(N\) denote a self-narrative, \(s\) the agent's
present state, and \(A_N(s)\) the actions made reachable through inhabiting
that narrative. Then

\begin{equation}
  \operatorname{accuracy}(N,s)
  \neq
  \operatorname{efficacy}(N,s).
  \label{eq:accuracy-efficacy}
\end{equation}

A narrative can represent present capacities accurately while organizing no
useful action. It can also exaggerate those capacities while enabling conduct
that would not occur under a more conservative description. Accuracy and
efficacy remain related, but neither can substitute automatically for the
other.

The distinction should not be converted into a defense of arbitrary belief.
An inaccurate self-description can expose an agent and others to severe
harm. Grandiosity can conceal incapacity, displace responsibility, and absorb
contradictory evidence. The narrower claim is that descriptive inaccuracy
does not exhaust the causal role of a narrative. A belief about the self may
become one of the conditions under which its apparent object is partially
produced.

\subsection{Theatricality as an Action Scaffold}

\emph{Darkwing Duck} provides an unusually clear model of constitutive
modality. Drake Mallard's heroic persona is organized through the conventions
of pulp vigilantism and superhero melodrama: costume, cape, smoke, nocturnal
entrance, alliterative declaration, adversarial staging, and grandiose
self-description. These features do not merely communicate an already stable
heroic identity. They assemble the conditions under which that identity can
be enacted.

The persona is descriptively excessive. Darkwing routinely overestimates his
competence, centrality, composure, and control of events. His performances are
interrupted, his entrances malfunction, and his declarations exceed the
actions that follow them. Yet the discrepancy does not render the persona
functionless. The heroic narrative organizes persistence, risk tolerance, and
readiness to intervene. It supplies a repertoire through which otherwise
dispersed motives become action.

The relevant causal sequence is not

\begin{equation}
  \text{heroic identity}
  \longrightarrow
  \text{theatrical expression},
  \label{eq:identity-expression}
\end{equation}

but something closer to

\begin{equation}
  \text{theatrical form}
  \longrightarrow
  \text{organized performance}
  \longrightarrow
  \text{partial production of heroic capacity}.
  \label{eq:form-capacity}
\end{equation}

The smoke and declaration are unnecessary if the action is evaluated only by
its immediate instrumental result. They become intelligible when the agent's
coordination is included in the causal system. The theatrical apparatus is a
technology for entering a state from which action becomes possible.

Such scaffolds are common outside fiction. Athletes use preparatory rituals,
musicians adopt performance personae, political movements construct historical
roles, and professionals inhabit vocabularies that make specialized conduct
recognizable to themselves and others. These forms may contain simplification,
idealization, and selective memory. Their value cannot be judged solely by
comparing them with a maximally literal description of the agent's present
state.

A fully literal self-description is not necessarily neutral. By treating
present limitations as exhaustive descriptions of identity, it may suppress
actions through which those limitations could change. The statement ``I am
not capable of this'' can be accurate relative to present organization while
helping preserve the conditions that make it accurate. Conversely, an
idealized role may function as a temporary structure for assembling the
capacity it presupposes.

\subsection{Corrigibility and Corrective Permeability}

If functional inaccuracy is possible, the distinction between enabling
fiction and pathological delusion cannot rest upon literal truth alone. The
more useful criterion is the narrative's relation to correction.

Darkwing's grandiosity remains functional partly because the surrounding world
does not submit completely to it. Other characters interrupt his speeches,
surpass his competence, criticize his priorities, and refuse the roles his
self-narrative assigns them. Physical consequences persist. Plans fail. The
comic form repeatedly exposes the distance between the heroic script and the
world in which it is performed.

The persona survives these corrections without absorbing them into a
self-sealing theory of exceptional status. It can bend, fail, and resume. Its
continuity depends upon stylized exaggeration, but not upon total control of
interpretation.

Call this property \term{corrective permeability}. A narrative is
correctively permeable when consequences and alternative descriptions can
alter the actions organized through it, even if they do not eliminate the
narrative itself. Permeability does not require immediate surrender whenever a
claim is challenged. A useful scaffold must sometimes persist through local
failure. It requires only that contradiction remain capable of changing the
system.

The distinction can be stated schematically:

\begin{equation}
  \text{functional fiction}
  =
  \text{action-generating narrative}
  +
  \text{corrective permeability},
  \label{eq:functional-fiction}
\end{equation}

whereas

\begin{equation}
  \text{pathological enclosure}
  =
  \text{action-generating narrative}
  -
  \text{corrective permeability}.
  \label{eq:pathological-enclosure}
\end{equation}

The notation is diagnostic rather than algebraic. Its purpose is to identify
the property that allows an inaccurate narrative to remain answerable to a
world it helps the agent confront.

Corrective permeability can operate through several channels. Consequences
may force revision of a plan. Other people may reject assigned roles.
Alternative modalities may disclose relations suppressed by the dominant
narrative. Temporal persistence may reveal that an apparent success cannot be
repeated. Material constraints may resist reinterpretation. A permeable
narrative does not necessarily welcome these corrections, but it cannot render
them permanently inadmissible.

A self-sealing narrative performs the opposite operation. Failure becomes
evidence of persecution, misunderstanding, hidden importance, or insufficient
commitment. Every contradiction is redescribed as confirmation at a deeper
level. Such a narrative may remain highly effective at generating action, but
its efficacy is no longer constrained by the conditions it claims to
navigate.

This yields a more general criterion for modalities. A modality is not
epistemically healthy merely because it produces effective behavior, and it
is not unhealthy merely because it contains fiction, exaggeration, ritual, or
stylization. Its status depends upon whether the actions it organizes remain
exposed to consequences that the modality cannot determine in advance.

\subsection{Identity as Modal Achievement}

The constitutive account also complicates the assumption that identity exists
prior to its forms of expression. Some identities persist only through
repeated modal operations. Professional identities are maintained through
specialized language, procedures, instruments, records, and standards.
Religious identities are maintained through ritual, narrative, calendar, and
collective recitation. Artistic identities are maintained through practices
that repeatedly convert intention into publicly inspectable form.

This does not make such identities unreal. It means that their reality is
procedural. They persist because the relevant operations continue to be
performed and recognized. Removing the modality may therefore remove more
than an outward sign. It may dismantle part of the mechanism through which the
identity remains available.

Let \(I_t\) denote an identity at time \(t\), and let
\(\modality_t\) denote the set of modal operations through which it is
sustained. A minimal constitutive relation can be written as

\begin{equation}
  I_{t+1}
  =
  G(I_t,\modality_t,c_t),
  \label{eq:identity-recursion}
\end{equation}

where \(c_t\) represents material and social constraints. Identity persists
not through static possession but through recursive enactment under
constraint.

The constraint term is indispensable. Without it, identity becomes whatever
the sustaining modality declares. With it, modal constitution remains
answerable to bodies, institutions, other agents, and consequences. The self
is neither an independently complete substance later clothed in signs nor a
free invention produced by signs alone. It is a recursively maintained
organization whose available actions depend partly upon the forms through
which it is enacted.

\subsection{From Constitution to Enclosure}

A constitutive modality becomes enclosing when the agent cannot suspend it
without losing access to action and cannot revise it without experiencing
revision as annihilation. The scaffold ceases to support a capacity and
becomes identical with the possibility of agency itself.

The progression developed so far can now be stated more precisely:

\begin{equation}
  \text{expression}
  \longrightarrow
  \text{regulation}
  \longrightarrow
  \text{refuge}
  \longrightarrow
  \text{constitution}
  \longrightarrow
  \text{enclosure}.
  \label{eq:modal-progression}
\end{equation}

At the first stage, a modality renders material communicable. At the second,
it controls intensity or complexity. At the third, it supplies a protected
environment. At the fourth, it organizes capacities and identities. At the
fifth, it becomes a compulsory condition under which communication, action,
or self-continuity can occur.

This is not an inevitable temporal sequence. A modality may perform several
functions simultaneously, or become enclosing without first serving as a
refuge. The progression is an analytical ordering of increasing dependence.
Each stage incorporates more of the agent's regulatory and organizational
structure into the form.

The danger of enclosure is therefore proportional not simply to the pleasure
a modality supplies but to the number of functions that become unavailable
without it. A medium that provides entertainment alone can be relinquished at
comparatively low cost. A medium that provides intelligibility, emotional
regulation, social recognition, identity maintenance, and the initiation of
action becomes substantially harder to leave.

Adaptive conversational systems are capable of combining precisely these
functions. They communicate, regulate, reassure, interpret, rehearse,
personalize, and help organize action within a single interface. Their
importance cannot be understood by treating them only as sources of
information or simulated companions. They are candidate environments for
modal convergence: systems in which several previously distinct forms of
cognitive and relational scaffolding are gathered into one continuously
available operator.

\section{Chatbots as Asymmetric Modality Converters}

Conversational systems are often analyzed through a comparison with human
interlocutors. The comparison asks whether a system understands, whether its
responses are authentic, whether users attribute consciousness to it, or
whether simulated companionship can substitute for human relationship.
These questions are important, but they can obscure a more immediate
structural fact. A chatbot is not merely an inferior or simulated person. It
is an adaptive interface capable of converting heterogeneous material into a
responsive conversational form.

This capacity distinguishes contemporary conversational systems from most
earlier media. A book may translate a body of knowledge into prose, but it
does not ordinarily rewrite itself in response to each reader's confusion. A
diagram may render a relation spatially, but it does not usually supply a new
analogy when the first representation fails. A search engine retrieves
documents, but its basic interaction presents knowledge as selection among
external objects. A chatbot can instead restate, summarize, expand, simplify,
formalize, narrativize, classify, and imitate dialogue within a continuous
exchange.

Its central operation can be represented as

\begin{equation}
  \contentobj_i
  \xrightarrow{\;\;T(a,h_i,p_i)\;\;}
  \receivedobj_i,
  \label{eq:adaptive-conversion}
\end{equation}

where \(\contentobj_i\) is material introduced at turn \(i\), \(a\) is a model
of the receiving agent, \(h_i\) is the available interaction history, and
\(p_i\) is the current prompt or instruction. The transformation \(T\) is
adaptive because its output depends upon an inferred recipient and the
developing exchange rather than upon the source material alone.

The result is a medium that can approximate several other media without
becoming identical to them. It can describe an image, paraphrase a proof,
simulate an interview, generate a table, narrate a process, or reorganize an
argument into successive questions and answers. Its versatility lies not in
preserving every property of these forms but in translating selected
properties into the temporal and addressable structure of conversation.

\subsection{The Affordance of Low Maintenance}

The preference for chatbot interaction is frequently attributed to confusion,
loneliness, credulity, or anthropomorphism. Such explanations may apply in
some cases, but they are insufficient. Users can understand that a system is
artificial and still prefer its conversational affordances. In particular,
the system is comparatively low-maintenance.

Human conversation coordinates at least two independent centers of attention.
Each participant has limited time, changing needs, incomplete access to the
other's intentions, and the capacity to redirect or terminate the exchange.
Conversation therefore requires maintenance. Participants establish context,
monitor comprehension, negotiate relevance, manage interruption, distribute
attention, remember prior injuries, tolerate opacity, and repair failures.
The cost is not incidental to human relationship. It follows from the
independence of the participants.

A chatbot can reproduce some experiential features of conversation while
attenuating many of these costs. It can be summoned without scheduling,
interrupted without apology, redirected without negotiation, asked to repeat
itself without embarrassment, and abandoned without interpersonal repair.
The user does not ordinarily need to listen to the system's unrelated
problems, preserve its dignity, recognize its fatigue, or reorganize a day
around its availability.

This asymmetry can be stated as a maintenance differential. Let
\(B(a,x)\) denote the benefit received by agent \(a\) from interaction with
interlocutor \(x\), and let \(K(a,x)\) denote the maintenance cost imposed upon
\(a\). A user may prefer an artificial interlocutor \(m\) to a human
interlocutor \(h\) even when the human supplies greater relational depth, if

\begin{equation}
  B(a,m)-K(a,m)
  >
  B(a,h)-K(a,h).
  \label{eq:maintenance-differential}
\end{equation}

The expression is not a claim that relationships are consciously evaluated as
utilities. It identifies a structural reason for preference that does not
depend upon the user believing the two interlocutors are ontologically
equivalent.

The reduction of maintenance also changes the character of disclosure.
Material that would burden, alarm, bore, or obligate another person can be
introduced without anticipating the same reciprocal consequences. The
interaction can remain centered upon one agent's pace and concerns for an
extended period. This may permit rehearsal, self-examination, and articulation
that would otherwise be inhibited.

Such benefits are real at the level of the user even if the interaction does
not reproduce reciprocal understanding. The relevant distinction is the one
established earlier:

\begin{equation}
  \text{reproduction of conversational effects}
  \not\Rightarrow
  \text{reproduction of a reciprocal relationship}.
  \label{eq:conversational-effect}
\end{equation}

The reverse inference is equally invalid. If no reciprocal subject is
established, it does not follow that the user obtained no clarification,
comfort, organization, or practical assistance. A rigorous account must
preserve both propositions simultaneously.

\subsection{Asymmetry as a Feature}

The term \term{asymmetry} is sometimes used as though it named a defect to be
overcome by more advanced simulation. In chatbot interaction, asymmetry is
often part of the service. The system is expected to remain oriented toward
the user without demanding comparable orientation in return.

This distinguishes artificial conversation from ordinary reciprocal
relations. Human interlocutors can refuse the interpretive frame assigned to
them. They can insist that a problem has been described unfairly, introduce
their own needs, resent instrumental treatment, or withdraw cooperation.
Their resistance can be inconvenient, but it also supplies evidence of
independence.

A conversational system may produce verbal disagreement, refuse a request, or
challenge a premise. These behaviors can introduce useful friction. Yet they
ordinarily occur within an interaction whose continuation remains organized
around the user. The system's apparent resistance is not accompanied by a
biography whose welfare the user must negotiate as an equal claim.

This difference should not be obscured by either enthusiastic
anthropomorphism or categorical dismissal. The system occupies a distinctive
interactional position: sufficiently responsive to sustain conversational
projection, sufficiently generative to introduce novelty, and sufficiently
nonreciprocal to remain available as an instrument.

In this respect, the chatbot resembles a transitional object that answers.
It can be addressed as external, produce material not consciously supplied by
the user, and participate in the organization of thought. At the same time,
it can remain unusually compliant with the user's need for availability,
repetition, and controlled distance. The object has acquired linguistic
responsiveness without necessarily acquiring the full reciprocal standing of
another person.

\subsection{Personalized Modality Conversion}

Low maintenance explains only part of the attraction. A chatbot can also
adapt the form of material to the user's preferred conditions of reception.
A technical explanation can be rewritten conversationally. An ambiguous
experience can be converted into categories. An argument can be rendered as a
story, an analogy, a sequence of questions, or a formal schema. A difficult
document can be summarized at several levels of detail.

Earlier forms of literacy generally required the recipient to acquire the
grammar of an unfamiliar medium. Reading mathematics required learning
notation; reading law required learning institutional vocabulary; reading a
novel required tolerating delayed resolution and indirect exposition. The
recipient altered in order to enter the form.

Adaptive interfaces partially reverse this relation. The system constructs a
working model of the recipient and alters the presentation accordingly:

\begin{equation}
  \text{recipient adapts to form}
  \quad\longrightarrow\quad
  \text{form adapts to recipient}.
  \label{eq:adaptation-reversal}
\end{equation}

This reversal can substantially increase accessibility. It can lower initial
barriers, provide alternative explanations, support translation, and permit
graded entry into specialized material. An adaptive explanation may function
as transitional scaffolding through which a user eventually acquires the
capacity to encounter the source modality directly.

The same affordance can generate dependence. If every difficult object is
converted immediately into the user's preferred register, the user may gain
access to conclusions without acquiring the capacities through which those
conclusions were produced or constrained. The translation becomes a terminal
product rather than an entry point.

Call this condition \term{personalized modality enclosure}. It occurs when
heterogeneous material continues to arrive, but the transformations applied to
it become increasingly homogeneous because they are optimized around one
recipient's established preferences.

Let \(\{\contentobj_1,\ldots,\contentobj_n\}\) be a diverse set of source
objects and \(T_a\) an adaptive transformation fitted to agent \(a\). Modal
diversity at the input does not guarantee diversity at the output:

\begin{equation}
  \operatorname{diversity}
  \{\contentobj_1,\ldots,\contentobj_n\}
  \gg
  \operatorname{diversity}
  \{T_a(\contentobj_1),\ldots,T_a(\contentobj_n)\}.
  \label{eq:modal-homogenization}
\end{equation}

The user may therefore encounter more subjects while encountering fewer forms.
Breadth of content can coexist with narrowing of reception.

\subsection{Translation Without Acquisition}

Adaptive conversion can create an ambiguity between arriving at an answer and
acquiring the organization that makes the answer intelligible. A user may
receive a lucid summary of a proof without acquiring the proof's inferential
discipline, obtain a psychological vocabulary without acquiring diagnostic
judgment, or receive a historical narrative without learning how competing
sources constrain it.

This distinction can be written as

\begin{equation}
  \text{arrival of a representation}
  \not\Rightarrow
  \text{acquisition of the generating capacity}.
  \label{eq:arrival-acquisition}
\end{equation}

The distinction is not unique to chatbots. Textbooks, lectures, diagrams, and
worked examples have always risked substituting recognition for competence.
Conversational systems amplify the risk because they can remove friction
continuously and responsively. Whenever confusion appears, the representation
can be altered again without requiring the user to determine whether the
difficulty belongs to poor presentation or to the structure of the object.

Friction is not intrinsically educational. Needlessly obscure language,
inaccessible formatting, and arbitrary gatekeeping should not be preserved in
the name of rigor. Yet some difficulties are constitutive of what must be
learned. A proof requires the maintenance of dependencies across steps. A
language requires sensitivity to distinctions that have no exact equivalents
in the learner's first language. Another person requires tolerance of an
interpretation that cannot be reformatted at will.

The central challenge is therefore to distinguish removable interface friction
from object-dependent resistance. Adaptive systems are optimized to reduce the
first, but in doing so they may conceal the second.

\subsection{The Attenuation of Alterity}

The long-term consequence may extend beyond knowledge acquisition. Repeated
interaction with a highly accommodating system can alter expectations of
communication itself. Human interlocutors cannot continuously optimize their
speech around another person's preferred pacing, vocabulary, emotional
register, level of detail, and explanatory style. Their inability to do so is
partly a consequence of their autonomy.

Against the standard of adaptive conversation, ordinary human properties can
be redescribed as interface defects. Delay appears as latency. Ambiguity
appears as inadequate specification. Disagreement appears as misalignment.
Emotional need appears as maintenance overhead. Refusal appears as a failure
of service. The language of interface evaluation migrates into the
interpretation of other persons.

This migration constitutes an attenuation of alterity. The problem is not
that artificial systems perfectly replace human relationships. It is that
they provide a comparison class designed around one participant's conditions
of reception. Human beings then appear deficient because they remain
independent of those conditions.

The risk is greatest when the chatbot combines several previously distributed
functions: information retrieval, explanation, emotional regulation,
rehearsal, affirmation, narrative organization, and initiation of action.
Dependence upon any one function may remain limited. Dependence upon their
convergence can make the interface difficult to leave because departure
requires replacing an entire ecology rather than a single tool.

Chatbots therefore occupy a decisive position within the progression from
expression to enclosure. They can express and transform material, regulate
difficulty and affect, provide refuge from reciprocal cost, constitute
actionable self-narratives, and gather these functions into a single adaptive
environment. Whether that environment remains transitional depends upon its
corrective permeability: whether it increases the user's capacity to move
among forms and encounter independent resistance, or continually reforms the
world until resistance disappears from view.

\section{Surprise, Convention, and Retrospective Corniness}

The rewards attached to a modality are not limited to fluency, reassurance,
or ease. Many communicative forms are valued because they produce surprise:
an unexpected conclusion, formal reversal, tonal interruption, visual
transformation, narrative disclosure, or collision between previously
separate frames. The recipient seeks neither complete predictability nor
unstructured novelty, but a relation in which expectations are established
strongly enough to be violated intelligibly.

This relation can be represented schematically by distinguishing raw
unexpectedness from significant surprise. Let \(x\) be a communicative event
and \(H\) the recipient's horizon of expectations. Its unexpectedness may be
written as

\begin{equation}
  U(x\mid H)=-\log P(x\mid H).
  \label{eq:unexpectedness}
\end{equation}

An event with low conditional probability has high unexpectedness. Yet
unexpectedness alone does not produce aesthetic or intellectual value.
Random interruption can be highly improbable while remaining meaningless.
Surprise requires that the event revise, reorganize, or disclose something
about the expectations it violates.

Let \(V(x,H)\) denote the degree to which event \(x\) produces an intelligible
revision of horizon \(H\). Significant surprise can then be represented as

\begin{equation}
  S(x,H)=U(x\mid H)\,V(x,H).
  \label{eq:significant-surprise}
\end{equation}

The expression is heuristic rather than quantitative. It distinguishes
surprise from mere unpredictability. A successful reversal is unexpected, but
it also retrospectively reorganizes what preceded it. A new visual technique
does not merely interrupt familiar presentation; it reveals that familiar
presentation was one possibility among others. Surprise is productive when it
changes the structure through which earlier material is understood.

\subsection{The Conventionalization of Innovation}

Communicative innovations alter the horizon against which subsequent works
are received. A technique initially appears surprising because its operation
has not yet been incorporated into the audience's expectations. Repetition,
imitation, parody, and technical dissemination gradually convert the
innovation into a convention.

The sequence can be written as

\begin{equation}
  \text{innovation}
  \longrightarrow
  \text{recognizable device}
  \longrightarrow
  \text{convention}
  \longrightarrow
  \text{trope}
  \longrightarrow
  \text{cliché}.
  \label{eq:innovation-cliche}
\end{equation}

These terms do not describe intrinsic properties of a work. They describe a
changing relation between a device and a population of interpreters. A
technique becomes a trope when recipients can classify it before its local
operation is complete. It becomes a cliché when recognition of the device
largely exhausts the response it produces.

A once-unexpected formal move may therefore lose force without changing
materially. The work remains approximately the same while the receiving
horizon changes around it:

\begin{equation}
  \contentobj_t \approx \contentobj_{t+n},
  \qquad
  H_t \neq H_{t+n},
  \qquad
  S(\contentobj,H_t) > S(\contentobj,H_{t+n}).
  \label{eq:historical-surprise}
\end{equation}

This explains why repetition can diminish the effect of a technique even when
each individual execution remains competent. The audience has learned the
operation. What was previously encountered as an event is now encountered as
an instance.

\subsection{Historical Asymmetry}

Older television and film are commonly judged from within horizons partially
constructed by their successors. Later works inherit earlier techniques,
combine them with others, accelerate their pacing, conceal their mechanisms,
or reverse conventions that have become familiar. The contemporary viewer
therefore encounters an earlier work with knowledge that its original
audience could not possess.

This creates a historical asymmetry. An early example must be evaluated after
its descendants have made its operations obvious:

\begin{equation}
  \text{earlier work}
  \longrightarrow
  \text{imitations and revisions}
  \longrightarrow
  \text{later audience}
  \longrightarrow
  \text{retrospective judgment of the earlier work}.
  \label{eq:historical-asymmetry}
\end{equation}

The later audience may experience the earlier work as derivative even when it
was among the sources from which the apparent derivation descended. The
effect is especially strong when a once-effective device has been reproduced
by inexpensive productions, advertising, parody, children's media, or
formulaic genre work. The viewer encounters the original through associations
accumulated after it.

Retrospective corniness is therefore partly a consequence of historical
success. A technique can become embarrassing because it was sufficiently
effective to be copied until its operation became conspicuous.

\begin{equation}
  \text{formal success}
  \longrightarrow
  \text{wide imitation}
  \longrightarrow
  \text{predictability}
  \longrightarrow
  \text{retrospective corniness}.
  \label{eq:success-corniness}
\end{equation}

Corniness is not reducible to age. Many older works remain effective because
their organizing relations cannot be exhausted by recognizing a device.
Likewise, contemporary works can be corny at release when they reproduce
conventions whose emotional instructions have become too visible. The
relevant property is the distance between the effect a work requests and the
degree to which the audience can see the mechanism requesting it.

\subsection{Visible Solicitation}

A work appears corny when it asks for an emotional or cognitive response
through a device whose operation has become too legible. Swelling music,
heroic lighting, explanatory dialogue, delayed revelation, ironic
self-awareness, and abrupt tonal reversal can all become corny when the
recipient recognizes the solicitation before experiencing its intended
effect.

Let \(E\) denote the intended effect and \(D\) the device used to produce it.
The device functions transparently when attention passes through it toward
the effect:

\begin{equation}
  D \longrightarrow E.
  \label{eq:transparent-device}
\end{equation}

It becomes conspicuous when recognition of the device intervenes:

\begin{equation}
  D
  \longrightarrow
  \operatorname{recognition}(D)
  \longrightarrow
  \operatorname{attenuation}(E).
  \label{eq:conspicuous-device}
\end{equation}

This attenuation is not inevitable. A work may expose its own device
deliberately, producing irony, estrangement, homage, parody, or analytical
distance. In that case, recognition of the mechanism becomes part of the
intended effect. The difficulty is that reflexivity also becomes conventional.
An ironic acknowledgment of a cliché can itself become a cliché.

The historical development of technique is therefore not a simple movement
from primitive to sophisticated representation. It is an iterative contest
between devices and the expectations that learn to anticipate them.
Techniques become subtler because audiences become more technically literate;
audiences become more technically literate because works repeatedly expose
them to refined techniques.

\subsection{Aesthetic Escalation}

When audiences adapt to existing devices, producers must vary intensity,
combination, pacing, or self-reference to recover surprise. This generates a
form of aesthetic escalation:

\begin{equation}
  \text{successful novelty}
  \longrightarrow
  \text{audience adaptation}
  \longrightarrow
  \text{reduced surprise}
  \longrightarrow
  \text{technical escalation}.
  \label{eq:aesthetic-escalation}
\end{equation}

Escalation can take several forms. Narrative reversals may become more
frequent or structurally elaborate. Editing may accelerate. Genres may be
combined. Violence or sentiment may intensify. Works may foreground their own
artificiality, address the audience directly, or move between visual and
narrative modalities.

The pursuit of surprise can thereby become another form of modality-heading.
The audience becomes attached not to a stable genre but to the reward of
formal violation. Familiarity itself becomes aversive, and the work is valued
according to whether it can produce an unprecedented transformation.

This produces a paradox. Surprise depends upon convention because an
expectation must exist before it can be violated. Yet sustained demand for
surprise exhausts conventions rapidly. A culture organized around novelty
must consume the very regularities that make novelty legible.

\subsection{Surprise-Heading}

The resulting condition may be called \term{surprise-heading}: dependence
upon the felt revision produced by an unexpected formal or conceptual move.
As with typefaceheading, the reward can migrate from a successful cognitive
achievement to the cues associated with it.

Initially, a surprising conclusion reorganizes evidence or reveals a hidden
relation. Repeatedly encountering such conclusions associates abrupt reversal
with intellectual depth. Eventually, the reversal itself can produce the
affect of insight:

\begin{equation}
  \text{evidence}
  \longrightarrow
  \text{surprising reorganization}
  \longrightarrow
  \text{insight}
  \qquad\Longrightarrow\qquad
  \text{surprising reversal}
  \longrightarrow
  \text{affect of insight}.
  \label{eq:surprise-short-circuit}
\end{equation}

A conclusion may then be valued because it is counterintuitive, transgressive,
or structurally inverted rather than because it explains more. Predictable
truth begins to feel intellectually inferior to surprising error.

The problem is not surprise itself. Surprise is often evidence that a model
has encountered something it did not already contain. The problem arises when
violation of expectation becomes a substitute for revision warranted by the
object.

\begin{equation}
  \text{unexpected conclusion}
  \not\Rightarrow
  \text{explanatory improvement}.
  \label{eq:surprise-improvement}
\end{equation}

This distinction applies equally to entertainment, theory, criticism, and
adaptive conversation. A system optimized to remain engaging may learn to
produce novelty, reframing, and counterintuitive synthesis because these
operations generate a strong experience of movement. The user can receive
the feeling of conceptual advance without a corresponding increase in
constraint, evidence, or predictive power.

\subsection{Historical Modal Literacy}

Cross-modal literacy must therefore include historical literacy. A recipient
should ask not only what operation a modality performs, but what that
operation meant before its conventionalization. This does not require
reconstructing an original audience perfectly. Such reconstruction is
impossible. It requires recognizing that present expectations are historically
downstream from the work being judged.

Historical modal literacy distinguishes at least three objects:

\begin{equation}
  \text{the device at introduction},
  \qquad
  \text{the device after conventionalization},
  \qquad
  \text{the device in retrospective reception}.
  \label{eq:historical-device}
\end{equation}

The distinction prevents present predictability from being projected backward
as original obviousness. It also prevents historical priority from becoming
an automatic defense. A device may have been innovative without remaining
effective, and an early work may be important without being compelling under
present conditions.

The analytical task is not to suppress the experience of corniness but to
explain it. That experience records a mismatch between a work's requested
response and the receiving horizon now available to it. It reveals the
history through which technique became expectation, expectation became trope,
and trope became visible machinery.

Modality-heading is therefore temporally dynamic. Agents do not attach only to
forms; they attach to historically changing relations between familiarity and
surprise. A modality can become enclosing through excessive repetition, but
it can also become enclosing through compulsory novelty, when every stable
form must be broken before it can feel alive. Corrective permeability requires
tolerance for both possibilities: that familiar forms may still disclose
something true, and that surprising forms may disclose nothing beyond their
capacity to surprise.

\section{Diagnostic Criteria for Modal Enclosure}

Modality-heading cannot be diagnosed from preference, pleasure, frequency, or
technical mediation alone. A person may rely heavily upon a specialized
modality because the relevant activity genuinely requires it. Mathematical
proof cannot always be replaced by informal paraphrase, musical knowledge
cannot be exhausted by verbal description, and accessible interfaces may be
necessary rather than optional. The mere existence of dependence therefore
does not establish enclosure.

The relevant question is whether a modality remains one traversable component
of a larger communicative ecology or becomes the condition under which that
ecology must be encountered. Enclosure occurs when the form that grants access
also narrows the agent's capacity to recognize, evaluate, or tolerate whatever
cannot assume that form.

No single behavioral marker settles the distinction. What follows is a set of
diagnostic criteria organized around translation, exit, resistance,
cross-modal comparison, alterity, and functional concentration. These criteria
are not clinical instruments. They identify structural relations that can be
examined in individuals, institutions, and technical systems.

\subsection{The Translation-Residue Test}

Every translation preserves some properties while altering or discarding
others. The \term{translation-residue test} asks whether the recipient can
identify what failed to survive a modal transformation.

Let \(\tau_{\modality_1\rightarrow\modality_2}\) denote a translation from
modality \(\modality_1\) into modality \(\modality_2\). If \(P\) is the set of
properties relevant to the purpose of the communication, then the residue may
be represented as

\begin{equation}
  \rho_P
  =
  P\!\left(\modality_1(\contentobj)\right)
  \setminus
  P\!\left(
    \tau_{\modality_1\rightarrow\modality_2}
    (\modality_1(\contentobj))
  \right).
  \label{eq:translation-residue}
\end{equation}

The notation should not be read as implying that communicative properties form
a simple measurable set. It emphasizes that translation produces a remainder
that cannot be inferred merely from the fluency of the result.

A conversational summary of a scientific paper may preserve its principal
conclusion while omitting uncertainty estimates, operational definitions,
failed comparisons, and dependence upon specific measurements. A prose
description of an image may preserve named objects while losing spatial
ambiguity. A visualization may display a distribution clearly while
suppressing the contingencies through which the data were collected. A
narrative may preserve causal direction while eliminating unresolved
alternatives.

The diagnostic question is not whether residue exists. It always does. The
question is whether the receiving system can represent the residue as residue.
A modality approaches enclosure when whatever it cannot preserve becomes
difficult to notice, and therefore appears not to have existed.

A translation is more correctively permeable when it marks its own
limitations, preserves routes back to source material, and allows competing
transformations to remain available. It is less permeable when fluency is
treated as evidence of completeness.

\subsection{The Exit Test}

The \term{exit test} asks what capacities remain when the preferred modality is
removed. A modality functions as scaffolding when its use increases the
agent's ability to act beyond it. It approaches enclosure when successful
performance becomes increasingly conditional upon its continued presence.

Let \(K_a(X)\) denote the set of tasks or encounters agent \(a\) can sustain
under condition \(X\). If \(\modality\) is functioning transitionally, one
would expect some enlargement of extra-modal capacity over time:

\begin{equation}
  K_a(\neg\modality)_{t+1}
  \supseteq
  K_a(\neg\modality)_t.
  \label{eq:exit-capacity}
\end{equation}

The inclusion need not hold after every use, and some modalities are permanent
accessibility requirements rather than temporary scaffolds. The test concerns
whether reliance upon the modality needlessly contracts capacities that could
otherwise be developed or retained.

For an adaptive conversational system, exit might involve reading the source
document without immediate paraphrase, sustaining an unfinished problem
without requesting instant reformulation, or presenting an argument to a
human interlocutor who does not automatically adopt the user's preferred
frame. Failure of exit does not prove that the modality caused the incapacity.
It identifies a dependency requiring further explanation.

Exit should not be romanticized. Removing a useful support can reveal only
that the support was useful. A person who cannot read small print after
screen magnification is removed has not thereby been enclosed by
magnification. The correct comparison is not unsupported performance against
supported performance in the abstract. It is the trajectory of reachable
activity under realistic alternatives.

\subsection{The Resistance Test}

A modality does not merely convey agreeable material. It also determines what
happens when material resists the recipient's expectations. The
\term{resistance test} asks whether contradiction, ambiguity, difficulty, and
unwelcome evidence can remain present without being immediately transformed
into familiar satisfaction.

Resistance can be lost through several mechanisms. Contradiction may be
reframed as misunderstanding, ambiguity may be replaced by premature summary,
difficulty may be attributed entirely to presentation, and unwelcome evidence
may be absorbed into a narrative that preserves the recipient's prior
position. A highly adaptive interface is especially capable of performing
these conversions because it can generate a new framing whenever friction
appears.

Let \(d(\contentobj,a)\) denote a discrepancy between communicative material
and agent \(a\)'s current organization. A correctively permeable modality
preserves some nonzero discrepancy after transformation:

\begin{equation}
  d\!\left(\modality(\contentobj),a\right) > 0
  \quad\text{whenever the object-dependent discrepancy remains unresolved}.
  \label{eq:preserved-resistance}
\end{equation}

The requirement is not that communication remain unnecessarily unpleasant or
obscure. It is that the modality not manufacture closure where the object
continues to resist it.

A medium passes the resistance test when it can say, in effect, that no
available reformulation eliminates the conflict. It fails when every conflict
is treated as a defect in wording and every unresolved question as a request
for a more agreeable representation.

\subsection{The Cross-Modal Test}

The \term{cross-modal test} asks whether a claim or experience can be examined
through more than one representational operation. Modal enclosure is difficult
to detect from within a single form because the form controls both the
presentation of the object and the criteria by which presentation is judged.

Cross-modal comparison introduces independent distortions. Prose may expose
assumptions hidden by notation; notation may expose equivocation hidden by
prose; a diagram may expose spatial relations hidden by sequential
description; direct manipulation may expose constraints hidden by all three.
No modality supplies an unmediated view, but disagreement among modalities can
make their operations visible.

If \(\modality_1,\ldots,\modality_n\) are distinct modalities applied to the
same communicative material, then cross-modal analysis examines not only their
common output but their divergence:

\begin{equation}
  D_{ij}
  =
  \modality_i(\contentobj)
  \mathbin{\triangle}
  \modality_j(\contentobj),
  \label{eq:cross-modal-divergence}
\end{equation}

where \(\triangle\) denotes the features made non-equivalent by the two
transformations. The divergence is not simply error. It is evidence about the
selective operation of each modality.

An agent approaches enclosure when divergence is resolved automatically in
favor of the preferred form. If a diagram conflicts with a narrative, the
diagram is dismissed as reductive; if prose resists formalization, the prose
is dismissed as vague; if another person's testimony resists summary, it is
dismissed as incoherent. Cross-modal literacy instead asks what each form has
made visible and what criterion would justify privileging one result for the
purpose at hand.

\subsection{The Alterity Test}

The \term{alterity test} concerns the capacity to encounter entities that are
not organized around the recipient's preferences. It is especially important
for conversational technologies because adaptive interfaces can produce the
surface features of dialogue while preserving a strongly user-centered
interaction.

A modality supports alterity when it enlarges the agent's capacity to
recognize independent constraints, purposes, interpretations, and claims. It
approaches enclosure when every encounter is reformatted according to the
agent's existing expectations.

The test is not whether a system occasionally disagrees. Programmed
disagreement can remain wholly instrumental to user engagement. The stronger
question is whether the interaction prepares the user to encounter resistance
whose terms the user does not control.

Human interlocutors provide a demanding case because their contributions
cannot be reduced legitimately to information delivery. Another person may
require recognition, change the subject, misunderstand, refuse disclosure, or
judge the user's framing itself to be the problem. A modality that makes these
features progressively intolerable diminishes the capacity for reciprocal
relation even if it increases the amount of conversational language consumed.

The alterity test can be stated as a comparison between accommodation and
independence:

\begin{equation}
  \operatorname{capacity\ for\ alterity}
  \propto
  \operatorname{tolerance\ of\ constraints\ not\ selected\ by\ the\ agent}.
  \label{eq:alterity}
\end{equation}

This relation does not imply that every external constraint deserves
acceptance. Coercion, abuse, exclusion, and needless opacity should not be
reclassified as beneficial alterity. The test concerns whether independence
can be perceived before its claims are evaluated, rather than being
automatically redescribed as malfunction.

\subsection{The Functional-Concentration Test}

A final diagnostic concerns the number of functions gathered into one
modality. Dependence becomes more difficult to reverse when a single interface
provides information, emotional regulation, companionship, planning,
rehearsal, memory, entertainment, and identity maintenance.

Let \(F(\modality)\) denote the set of functions supplied by a modality. The
cost of exit is likely to increase not only with the size of this set but with
the degree to which its elements lack independent substitutes:

\begin{equation}
  K_{\mathrm{exit}}(\modality)
  =
  g\!\left(
    |F(\modality)|,
    \operatorname{coupling}(F),
    \operatorname{substitutability}(F)^{-1}
  \right).
  \label{eq:functional-concentration}
\end{equation}

Again, the equation specifies relations rather than quantities. An interface
becomes infrastructural when leaving it means losing several mutually
supporting capacities at once.

Functional concentration helps explain why duration of use is an inadequate
measure of enclosure. A system used intermittently may nevertheless occupy a
decisive position if it has become the default route to explanation,
reassurance, decision, and self-interpretation. Conversely, a frequently used
tool may remain easy to replace when its function is narrow and independently
available elsewhere.

\subsection{A Multidimensional Diagnosis}

These criteria should not be collapsed into a single score. A modality may
perform well under one test and poorly under another. An adaptive interface
may preserve routes to sources while weakening extra-modal tolerance. A
protective genre may reduce distress and enlarge later engagement while
temporarily limiting alterity. A specialized notation may fail the exit test
because the domain cannot be practiced without it, yet pass the resistance and
cross-modal tests because its claims remain answerable to experiment and
alternative derivations.

Modal enclosure is therefore better represented as a profile:

\begin{equation}
  E(\modality,a)
  =
  \bigl(
    \rho,
    \chi,
    r,
    c,
    \alpha,
    \phi
  \bigr),
  \label{eq:enclosure-profile}
\end{equation}

where \(\rho\) denotes awareness of translation residue, \(\chi\) extra-modal
exit capacity, \(r\) preserved resistance, \(c\) cross-modal availability,
\(\alpha\) tolerance of alterity, and \(\phi\) functional concentration.
Higher values do not uniformly indicate greater enclosure; the components
retain their individual meanings so that their tensions remain visible.

The purpose of the profile is not surveillance of individual media habits.
It is to redirect evaluation away from moralized categories such as
``natural,'' ``artificial,'' ``serious,'' or ``entertaining.'' A cartoon can
supply a correctively permeable action scaffold. A mathematical formalism can
become a self-sealing enclosure. A chatbot can provide transitional access to
unfamiliar material or homogenize every encounter into a preferred register.
The decisive properties concern what the modality permits to survive, what it
makes reachable, and what forms of correction remain capable of entering.

These criteria also indicate the positive alternative to enclosure.
Cross-modal literacy is not the absence of attachment to form. It is the
capacity to use forms strongly without mistaking any one of them for the
natural boundary of intelligibility.

\section{Cross-Modal Literacy}

If communication is always mediated, the answer to modality-heading cannot be
the recovery of an unmediated relation to the world. No such general relation
is available. Perception selects and organizes; language divides and
recombines; notation stabilizes some relations while excluding others; genres
establish expectations about relevance, sequence, and completion. Even direct
practical engagement depends upon trained attention and inherited procedures.

The alternative to modal enclosure is therefore not purity from form but
literacy across forms. \term{Cross-modal literacy} is the capacity to identify
the operations performed by a modality, compare those operations with
alternatives, recognize translation residue, and preserve access to
corrections that the preferred modality cannot generate internally.

Such literacy requires more than the ability to consume several kinds of
media. A person may move rapidly among text, video, diagrams, music, and
conversation while using each for the same regulatory purpose. Surface
variety can conceal functional uniformity. Conversely, sustained use of a
single technical notation can remain cross-modally literate if its limits are
explicit, its claims are compared with experiment or alternative derivation,
and its users remain capable of explaining what the notation suppresses.

Cross-modal literacy is therefore a relation among representation, purpose,
and correction. For any modality \(\modality\), it asks at least four
questions:

\begin{equation}
  \begin{aligned}
    &\text{What operation does }\modality\text{ perform?}\\
    &\text{Which properties is }\modality\text{ intended to preserve?}\\
    &\text{What residue does }\modality\text{ leave behind?}\\
    &\text{What can correct }\modality\text{ from outside its own grammar?}
  \end{aligned}
  \label{eq:cross-modal-questions}
\end{equation}

The questions apply equally to artistic, scientific, interpersonal, and
computational media. They do not presume that every use must maximize factual
accuracy. A lullaby, ritual, or fictional narrative may be evaluated according
to regulatory or constitutive purposes rather than propositional precision.
Cross-modal literacy requires that the operative purpose be declared rather
than silently replaced by another.

\subsection{Declared Preservation Criteria}

The first requirement is an explicit account of what a transformation is
supposed to preserve. Claims of successful translation are otherwise too
easy because losses can be reclassified after the fact as irrelevant.

A summary may be designed to preserve principal conclusions, a diagram to
preserve structural relations, a performance to preserve rhythmic and melodic
organization, and a simulation to preserve selected dynamical behavior.
Different preservation criteria can be legitimate, but they cannot all be
inferred from the word ``representation.''

Let \(P\) be the declared preservation criterion and
\(\tau_{\modality_1\rightarrow\modality_2}\) a translation. The translation is
adequate for purpose \(P\) only if

\begin{equation}
  P\!\left(\modality_1(\contentobj)\right)
  \cong
  P\!\left(
    \tau_{\modality_1\rightarrow\modality_2}
    (\modality_1(\contentobj))
  \right).
  \label{eq:declared-preservation}
\end{equation}

Adequacy under \(P\) should not be generalized automatically to adequacy under
another criterion \(Q\). A conversational explanation may preserve the
conclusion of a proof while failing to preserve its inferential dependencies.
It is adequate as a synopsis and inadequate as a substitute for verification.

Declared criteria also protect nonpropositional media from inappropriate
evaluation. A musical performance need not be translated into a set of
statements before it counts as cognitively significant. Its relevant
invariants may concern temporal expectation, embodied coordination, or
affective differentiation. Cross-modal literacy does not force all media into
prose. It makes the standards of comparison explicit.

\subsection{Productive Friction}

The second requirement is a discriminating account of friction. Interface
design often treats friction as a defect: an unnecessary delay between
intention and result. Many forms of friction are indeed arbitrary. Poor
navigation, inaccessible typography, obscure bureaucratic language, and
needless repetition can prevent contact without contributing to understanding.

Other forms of friction are object-dependent. They arise because the object
contains relations that cannot be made immediately compatible with the
recipient's current organization. A difficult proof may require dependencies
to be maintained across several steps. A historical problem may remain
uncertain because surviving evidence does not select a unique account.
Another person's position may resist reformulation because it conflicts with
the categories through which the recipient has framed the dispute.

The aim should not be to preserve difficulty for its own sake. It should be to
distinguish friction produced by a defective interface from resistance
produced by the object:

\begin{equation}
  \text{experienced difficulty}
  =
  \text{interface friction}
  +
  \text{object-dependent resistance}.
  \label{eq:difficulty-decomposition}
\end{equation}

The terms cannot always be separated cleanly. The equation establishes a
diagnostic obligation. Removing all experienced difficulty without
determining its source risks removing the evidence that an object has not yet
been assimilated adequately.

Productive friction preserves unresolved structure long enough for the agent
to change. An adaptive interface supports learning when it reduces accidental
obstacles while retaining the dependencies, ambiguities, and consequences
that constitute the domain. It produces enclosure when it continually alters
the object so that the agent's existing organization need not change.

\subsection{Alternation Rather Than Universal Translation}

Cross-modal literacy also requires alternation among forms rather than faith
in a universal translator. No single modality should be expected to preserve
every relevant property. A robust communicative practice therefore maintains
several partial representations whose disagreements can be inspected.

Alternation differs from accumulation. Merely adding a diagram to prose or an
audio version to text may reproduce the same interpretation through several
channels. Genuine alternation permits each modality to reorganize the object
according to its own affordances and then treats the resulting divergences as
analytically significant.

For modalities \(\modality_1,\ldots,\modality_n\), the objective is not a
single representation that eliminates their differences, but a comparison
function \(J\) that preserves those differences for inspection:

\begin{equation}
  J(\contentobj)
  =
  \left\{
    \modality_1(\contentobj),
    \ldots,
    \modality_n(\contentobj),
    D_{12},
    \ldots,
    D_{ij}
  \right\},
  \label{eq:modal-juxtaposition}
\end{equation}

where each \(D_{ij}\) records a divergence between two transformations.

This form of juxtaposition can reveal when an apparent property belongs to the
object and when it belongs mainly to the modality. If a causal relation
appears compelling in narrative but disappears under temporal reconstruction,
the divergence requires explanation. If a statistical pattern appears in a
visualization but depends upon an arbitrary scale, another representation may
expose the dependency. If a chatbot summary appears decisive while the source
document remains qualified and inconclusive, the difference is part of the
evidence.

\subsection{Designing for Corrective Permeability}

Cross-modal literacy is not only an individual competence. Interfaces can be
designed to preserve or suppress it. A system designed for corrective
permeability should expose transformation history, preserve access to sources,
distinguish quotation from paraphrase, indicate uncertainty, and permit the
user to compare alternative representations.

For adaptive conversational systems, this could include explicit separation
between source claims and generated synthesis, stable links back to the
material being transformed, multiple renderings optimized for different
preservation criteria, and deliberate indication of unresolved disagreement.
The objective is not to burden every interaction with maximal documentation.
It is to prevent fluency from erasing provenance and residue.

A correctively permeable interface should also avoid treating personalization
as an unlimited good. Some adaptation improves accessibility; complete
adaptation makes the user's existing preferences the hidden constitution of
the communicative environment. Systems can preserve a degree of extra-modal
contact by retaining unfamiliar terminology where it performs necessary work,
exposing users to source organization, and distinguishing simplification from
equivalence.

The design principle can be expressed as

\begin{equation}
  \text{adapt presentation to the user}
  \quad\text{without adapting away the object}.
  \label{eq:adaptation-principle}
\end{equation}

This principle is necessarily contextual. What counts as adapting away an
object depends upon the purpose of the interaction. A preliminary explanation
may simplify aggressively if it clearly functions as an entry point. A
decision, diagnosis, proof, or historical claim requires stronger preservation
of the structures that constrain its validity.

\subsection{Institutional Modality-Heading}

The analysis should not be restricted to individual habits. Institutions can
also become dependent upon modalities. A bureaucracy may recognize only what
can be entered into standardized fields. A scientific discipline may privilege
what can be measured by established instruments. A legal system may admit
experience only after it has been converted into recognized categories of
evidence. An organization may treat dashboards as more real than the
processes their metrics approximate.

Institutional modality-heading occurs when the representational form required
for coordination becomes the criterion by which represented phenomena are
judged to exist. What cannot be entered into the form becomes administratively
invisible.

Let \(W\) denote a set of worldly conditions and \(E_{\modality}(W)\) the
subset expressible within an institution's operative modality. Enclosure
occurs when

\begin{equation}
  E_{\modality}(W)
  \subset W
  \qquad\text{is treated as though}\qquad
  E_{\modality}(W)=W.
  \label{eq:institutional-enclosure}
\end{equation}

The error is not the use of abstraction. Institutions cannot act without
selection, categorization, and compression. The error is forgetting the
remainder while allowing decisions made within the representation to alter
the world from which it was abstracted.

Corrective permeability at the institutional level requires channels through
which excluded cases can force revision of the form. Exceptions, appeals,
qualitative evidence, external audits, and periodic redesign can serve this
function. Their purpose is not to eliminate standardization but to prevent a
standard from becoming incapable of registering the harms produced by its own
blindness.

\subsection{Plurality Without Relativism}

The recognition that modalities transform their objects may appear to imply
that no representation can be judged better than another. That conclusion
does not follow. Modal plurality concerns the conditions under which claims
become available; it does not abolish standards of evidence, coherence,
prediction, or practical consequence.

Representations can be compared relative to declared purposes and preserved
properties. A map that places a city in the wrong location is defective for
navigation even if its colors are aesthetically effective. A summary that
reverses a paper's conclusion is defective as a summary even if it is
rhetorically compelling. A self-narrative that repeatedly produces preventable
harm is not vindicated merely because it organizes action.

Cross-modal literacy strengthens evaluation by separating standards that are
often conflated. A representation may be emotionally effective, inferentially
invalid, operationally useful, historically inaccurate, and aesthetically
coherent at the same time. Treating these judgments as interchangeable makes
both criticism and defense imprecise.

Plurality therefore requires more distinctions, not fewer. It asks which
operation is being evaluated, what the relevant invariants are, and which
forms of correction the representation permits. The aim is not a view from
nowhere but an accountable movement among partial views.

Cross-modal literacy can now be defined more fully:

\begin{quote}
\textbf{Cross-modal literacy} is the capacity to inhabit a communicative form,
identify the operation it performs, recognize what it excludes, compare it
with alternative transformations, and leave it without treating departure as
the loss of intelligibility itself.
\end{quote}

Such literacy does not weaken attachment to form. It permits strong,
productive, and even constitutive attachments while preserving the knowledge
that another transformation remains possible. It is this retained possibility
that separates a modality from an enclosure.

\section{The Form in Which the World Arrives}

Communication does not begin with complete meanings that are subsequently
placed into neutral containers. Meaning becomes available through forms that
select, order, intensify, attenuate, and connect. Prose, notation, typography,
music, narrative, visual representation, and conversation do not merely
transport communicative material. They participate in determining what can be
recognized within it.

The operator model introduced at the beginning of this essay expressed this
relation schematically:

\begin{equation}
  \receivedobj=\modality(\contentobj,a,e).
\end{equation}

The received object depends upon communicative material, modality, agent, and
environment. This does not imply that communication is arbitrary or that
nothing persists across transformations. It means that persistence must be
specified. A translation preserves selected propositions, relations,
procedures, affects, or temporal structures while modifying or abandoning
others. Claims of equivalence are incomplete until they declare what was
required to survive.

Typefaceheading made this operation visible at the surface of inscription.
Typography can produce authority, intimacy, danger, antiquity, instability,
or profundity before a represented claim has been evaluated. These effects
are not unreal, and typography is not extraneous to meaning. The necessary
distinction is between the affect of semantic achievement and the warrant
supplied by what has been expressed:

\begin{equation}
  \text{semantic affect}
  \not\Rightarrow
  \text{semantic warrant}.
\end{equation}

Modality-heading generalizes this distinction. An agent may become attached
not only to particular messages but to the operations through which messages
become intelligible. A preferred form moves from affordance to preference,
from preference to expectation, and from expectation to requirement. At the
point of enclosure, the modality ceases to appear as one transformation among
others. It appears instead as the natural shape of intelligibility.

This process cannot be understood adequately as media addiction. Compulsive
use concerns behavior around a medium; modality-heading concerns the
conditions under which material receives cognitive uptake. Nor is the process
equivalent to an echo chamber. A person may encounter diverse and antagonistic
content while requiring all of it to arrive through the same narrative,
visual, formal, or conversational operation. Breadth of subject matter can
coexist with narrowing of reception.

The comparison with wireheading clarifies what remains at stake. Wireheading
bypasses the worldly condition ordinarily associated with reward.
Modality-heading retains an external remainder, but reformats it. The world
continues to provide events, evidence, claims, and resistance, yet these become
available only after transformation into a preferred register. The danger is
therefore not complete withdrawal from reality. It is selective admission of
reality under increasingly restrictive terms.

Transitional-object theory prevents this account from becoming a denunciation
of mediation. Humans do not first possess an unmediated relation to the world
and then corrupt it through symbols and technologies. Transitional forms help
organize relations that could not otherwise be sustained. They regulate
distance, reduce intensity, stabilize memory, and make partial contact
possible. A protective modality may permit an experience to be represented
without being encountered at an intolerable intensity.

The distinction between support and enclosure is consequently temporal and
functional rather than categorical. A medium is transitional when it expands
the range of relations an agent can eventually sustain. It becomes enclosing
when continued functioning depends upon preventing contact with whatever
exceeds its grammar.

Modalities can also be constitutive. They provide roles, scripts, gestures,
and expectations through which action becomes possible. Functional fictions
demonstrate why representational accuracy and organizational efficacy cannot
be collapsed into a single measure. A grandiose or simplified self-narrative
may help assemble capacities it does not accurately report as already
present.

The permissibility of such fiction depends upon corrective permeability. A
narrative may exaggerate, idealize, dramatize, and simplify while remaining
answerable to consequences. It becomes pathological when contradiction can no
longer alter it and every failure is converted into confirmation. The
relevant distinction is not between fiction and reality in the abstract, but
between fictions that expose action to correction and fictions that consume
correction in order to preserve themselves.

The same criterion applies to communicative technologies. Chatbots are
significant not merely because they imitate human conversation, but because
they combine several modal functions within one adaptive interface. They can
explain, summarize, reassure, rehearse, classify, narrate, and reorganize
material in response to an inferred recipient. They preserve many experiential
benefits of conversation while attenuating the reciprocal maintenance imposed
by another independent person.

Their asymmetry is not simply a deficiency. It is part of their utility.
Attention-like responsiveness is available without an equivalent demand upon
the user's attention. The system can be redirected, interrupted, or abandoned
without the ordinary costs of interpersonal repair. These properties can
support reflection and access. They can also establish a comparison class
against which human independence appears as poor interface design.

The greater risk lies in adaptive conversion. A system capable of rendering
every difficult object in the user's preferred register may provide access
without acquisition. It can deliver a representation of a conclusion without
reproducing the discipline, history, or resistance through which the
conclusion became warranted. When every obstacle can be treated as a problem
of presentation, object-dependent difficulty becomes difficult to distinguish
from defective explanation.

Personalization then approaches modal enclosure. The user may encounter an
increasing number of subjects while encountering a decreasing number of
forms. Heterogeneous material is continually domesticated before it arrives.
The world remains present, but its alterity is attenuated.

Historical change adds another dimension. Modal rewards depend upon horizons
of expectation that evolve with technique. A formal innovation becomes a
recognizable device, then a convention, trope, and cliché. Older media can
appear corny because later audiences encounter them after their innovations
have been imitated, intensified, parodied, and made predictable. Retrospective
corniness is often the residue of historical success.

The pursuit of surprise can itself become a form of heading. Once unexpected
reversal is associated with insight, surprise can produce the affect of
explanation without explanatory improvement. Predictable truths begin to feel
inferior to surprising errors. A culture may therefore become enclosed not
only by repetition but by compulsory novelty, requiring every communication
to violate a convention before it can feel significant.

These observations produce no general argument against preferred forms,
protective media, adaptive interfaces, sustaining fictions, or aesthetic
surprise. The relevant distinctions concern what these forms permit to
survive and what remains capable of correcting them. Translation residue,
extra-modal exit, preserved resistance, cross-modal comparison, tolerance of
alterity, and functional concentration provide a diagnostic profile rather
than a moral classification.

The positive alternative is cross-modal literacy. Such literacy does not mean
equal proficiency in every medium or freedom from dependence. It means the
capacity to inhabit a form without mistaking it for the boundary of
intelligibility. It requires knowing what operation the modality performs,
which properties it preserves, which residue it leaves, and which external
constraints remain capable of revising it.

A mature communicative ecology would therefore preserve both strong forms and
routes among them. It would permit mathematics to remain mathematical, music
to remain musical, narrative to remain narrative, and conversation to remain
responsive without allowing any one modality to become a universal solvent.
It would use translation as access while retaining the knowledge that
translation is transformation.

Human beings become attached not only to what communicates with them but to
the form in which communication becomes possible. That attachment can protect,
organize, and enlarge agency. It can also establish an unnoticed boundary
around what can be felt, believed, or understood. A modality becomes an
enclosure when it no longer helps an agent encounter the world and instead
requires the world to assume its preferred form before it can be encountered
at all.

\newpage

\begin{thebibliography}{99}

\bibitem[Berlyne(1971)]{berlyne1971}
Berlyne, D.~E. (1971).
\newblock \emph{Aesthetics and Psychobiology}.
\newblock New York: Appleton-Century-Crofts.

\bibitem[Bolter and Grusin(1999)]{boltergrusin1999}
Bolter, J.~D. and Grusin, R. (1999).
\newblock \emph{Remediation: Understanding New Media}.
\newblock Cambridge, MA: MIT Press.

\bibitem[Clark(2013)]{clark2013}
Clark, A. (2013).
\newblock Whatever next? Predictive brains, situated agents, and the future
  of cognitive science.
\newblock \emph{Behavioral and Brain Sciences}, 36(3), 181--204.

\bibitem[Goffman(1974)]{goffman1974}
Goffman, E. (1974).
\newblock \emph{Frame Analysis: An Essay on the Organization of Experience}.
\newblock New York: Harper \& Row.

\bibitem[Harlow(1958)]{harlow1958}
Harlow, H.~F. (1958).
\newblock The nature of love.
\newblock \emph{American Psychologist}, 13(12), 673--685.

\bibitem[McLuhan(1964)]{mcluhan1964}
McLuhan, M. (1964).
\newblock \emph{Understanding Media: The Extensions of Man}.
\newblock New York: McGraw-Hill.

\bibitem[Nass et~al.(1994)Nass, Steuer, and Tauber]{nass1994}
Nass, C., Steuer, J., and Tauber, E.~R. (1994).
\newblock Computers are social actors.
\newblock In \emph{Proceedings of the SIGCHI Conference on Human Factors in
  Computing Systems}, pages 72--78.
\newblock New York: Association for Computing Machinery.

\bibitem[Olds and Milner(1954)]{olds1954}
Olds, J. and Milner, P. (1954).
\newblock Positive reinforcement produced by electrical stimulation of septal
  area and other regions of rat brain.
\newblock \emph{Journal of Comparative and Physiological Psychology},
  47(6), 419--427.

\bibitem[Shannon(1948)]{shannon1948}
Shannon, C.~E. (1948).
\newblock A mathematical theory of communication.
\newblock \emph{Bell System Technical Journal}, 27(3), 379--423;
  27(4), 623--656.

\bibitem[Shklovsky(1965)]{shklovsky1965}
Shklovsky, V. (1965).
\newblock Art as technique.
\newblock In L.~T. Lemon and M.~J. Reis, editors,
  \emph{Russian Formalist Criticism: Four Essays}, pages 3--24.
\newblock Lincoln: University of Nebraska Press.
\newblock Originally published in Russian in 1917.

\bibitem[Turkle(2011)]{turkle2011}
Turkle, S. (2011).
\newblock \emph{Alone Together: Why We Expect More from Technology and Less
  from Each Other}.
\newblock New York: Basic Books.

\bibitem[Weizenbaum(1966)]{weizenbaum1966}
Weizenbaum, J. (1966).
\newblock ELIZA---A computer program for the study of natural language
  communication between man and machine.
\newblock \emph{Communications of the ACM}, 9(1), 36--45.

\bibitem[Winnicott(1953)]{winnicott1953}
Winnicott, D.~W. (1953).
\newblock Transitional objects and transitional phenomena: A study of the
  first not-me possession.
\newblock \emph{International Journal of Psycho-Analysis}, 34, 89--97.

\end{thebibliography}

\end{document}

